\Article{20}
Having completely solved our problem for moduli of the form $8n+1$, we turn to the other case, where $p$ has the form $8n+5$. We can treat it more briefly, since all the arguments differ little from the preceding ones.

Since, for such a modulus, $-1$ belongs to class $C$, the numbers in the collections $A,B,C,D$ have their complements to the sum $p$ in the classes $C,D,A,B$, respectively. It follows readily that

\SourceLeaf{95}%
\begin{center}
\begin{tabular}{c|l}
symbol&\shortstack{denotes the number of\\solutions of the congruence}\\\hline
$(00)$&$1+\alpha+\gamma\equiv0$\\
$(01)$&$1+\alpha+\delta\equiv0$\\
$(02)$&$1+\alpha+\alpha'\equiv0$\\
$(03)$&$1+\alpha+\beta\equiv0$\\
$(10)$&$1+\beta+\gamma\equiv0$\\
$(11)$&$1+\beta+\delta\equiv0$\\
$(12)$&$1+\beta+\alpha\equiv0$\\
$(13)$&$1+\beta+\beta'\equiv0$\\
$(20)$&$1+\gamma+\gamma'\equiv0$\\
$(21)$&$1+\gamma+\delta\equiv0$\\
$(22)$&$1+\gamma+\alpha\equiv0$\\
$(23)$&$1+\gamma+\beta\equiv0$\\
$(30)$&$1+\delta+\gamma\equiv0$\\
$(31)$&$1+\delta+\delta'\equiv0$\\
$(32)$&$1+\delta+\alpha\equiv0$\\
$(33)$&$1+\delta+\beta\equiv0$
\end{tabular}
\end{center}
from which we immediately obtain six equations:
\[
\begin{gathered}
(00)=(22),\quad(01)=(32),\quad(03)=(12),\\
(10)=(23),\quad(11)=(33),\quad(21)=(30)
\end{gathered}
\]

Multiplying the congruence $1+\alpha+\gamma\equiv0$ by a number $\gamma'$ from the collection $C$, chosen so that $\gamma\gamma'\equiv1$, and taking for $\gamma''$ the least residue of the product $\alpha\gamma'$, which must evidently also belong to the collection $C$, gives $\gamma'+\gamma''+1\equiv0$, from which we conclude that $(00)=(20)$.

In exactly the same way we obtain the equations $(01)=(13)$, $(03)=(31)$, $(10)=(11)=(21)$.

With these eleven equations we can reduce our sixteen unknowns to five and write the scheme $S$ as follows:
\[
\begin{array}{llll}
h,&i,&k,&l\\
m,&m,&l,&i\\
h,&m,&h,&m\\
m,&l,&i,&m
\end{array}
\]


\SourceLeaf{96}%
We also have the equations
\[
\begin{aligned}
(00)+(01)+(02)+(03)&=2n+1\\
(10)+(11)+(12)+(13)&=2n+1\\
(20)+(21)+(22)+(23)&=2n\\
(30)+(31)+(32)+(33)&=2n+1
\end{aligned}
\]
or, using the notation just introduced, these three (I):
\[
\begin{aligned}
h+i+k+l&=2n+1\\
2m+i+l&=2n+1\\
h+m&=n
\end{aligned}
\]
With their help we can therefore now reduce our unknowns to two.

We shall derive the remaining equations by considering the number of solutions of the congruence $1+\alpha+\beta+\gamma\equiv0$ (here too, $\alpha,\beta,\gamma$ denote unspecified numbers from the collections $A,B,C$, respectively). Consider \emph{first} that $1+\alpha$ gives $h,i,k,l$ numbers belonging to $A,B,C,D$, respectively, and that for each given value of $\alpha$ in these four cases there are $m,l,i,m$ solutions, respectively. The total number of solutions will be
\[
=hm+il+ik+lm
\]
\emph{Second}, since $1+\beta$ gives $m,m,l,i$ numbers belonging to $A,B,C,D$, and for each \emph{given} value of $\beta$ in these four cases there are $h,m,h,m$ solutions, the total number of solutions will be
\[
=hm+mm+hl+im
\]
From this we derive the equation
\[
0=mm+hl+im-il-ik-lm
\]
which, using the equation $k=2m-h$ obtained from (I), becomes
\[
0=mm+hl+hi-il-im-lm
\]
Now from equations I we also have $l+i=1+2h$, and hence
\[
\begin{aligned}
2i&=1+2h+(i-l)\\
2l&=1+2h-(i-l)
\end{aligned}
\]

\SourceLeaf{97}%
Substituting these values into the preceding equation gives
\[
0=4mm-4m-1-8hm+4hh+(i-l)^2
\]
If, finally, we replace $4m$ here with $2(h+m)-2(h-m)$, or, by the last equation in I, with $2n-2(h-m)$, we obtain
\[
0=4(h-m)^2-2n+2(h-m)-1+(i-l)^2
\]
and hence
\[
8n+5=(4(h-m)+1)^2+4(i-l)^2
\]
Thus, putting
\[
4(h-m)+1=a,\quad2i-2l=b
\]
gives
\[
p=aa+bb
\]

Since here too $p$ can be split in only one way into two squares, one even and the other odd, $aa$ and $bb$ will be completely determined numbers: evidently $aa$ must be set equal to the odd square and $bb$ to the even one. Moreover, the \emph{sign} of $a$ must be chosen so that $a\equiv1\ (\mathrm{mod.}\ 4)$, and the sign of $b$ so that $b\equiv af\ (\mathrm{mod.}\ p)$, as is readily proved by arguments entirely similar to those used in the preceding article.

With these preliminaries, the five numbers $h,i,k,l,m$ are determined in terms of $a,b$ and $n$ as follows:
\[
\begin{aligned}
8h&=4n+a-1\\
8i&=4n+a+2b+3\\
8k&=4n-3a+3\\
8l&=4n+a-2b+3\\
8m&=4n-a+1
\end{aligned}
\]
Or, if we prefer expressions in terms of $p$, the terms of the scheme $S$, multiplied by $16$, will be as follows:
\[
\begin{array}{l|l|l|l}
p+2a-7&p+2a+4b+1&p-2a+1&p+2a-4b+1\\
p-2a-3&p-2a-3&p+2a-4b+1&p+2a+4b+1\\
p+2a-7&p-2a-3&p+2a-7&p-2a-3\\
p-2a-3&p+2a-4b+1&p+2a+4b+1&p-2a-3
\end{array}
\]

\clearpage
\SourceLeaf{98}%
\Article{21}
Having solved our problem, we return to the main investigation, now undertaking the complete determination of the collection to which the number $2$ belongs.

I. Whenever $p$ has the form $8n+1$, we already know that the number $2$ lies either in the collection $A$ or in the collection $C$. In the former case we readily see that the numbers $\frac12(p-1)$, $\frac12(p+1)$ also belong to $A$, and in the latter to $C$. Now consider that, if $\alpha$ and $\alpha+1$ are consecutive numbers in the collection $A$, then $p-\alpha-1$, $p-\alpha$ are also such numbers; or, equivalently, numbers in the collection $A$ that are followed by a number in the same collection are always paired ($\alpha$ and $p-1-\alpha$). Their number, $(00)$, will therefore always be even, unless one is paired with itself, that is, unless $\frac12(p-1)$ belongs to $A$, in which case their number will be odd. From this we conclude that $(00)$ is odd whenever $2$ belongs to the collection $A$, and even whenever $2$ belongs to $C$. But we have
\[
16(00)=aa+bb-6a-11
\]
or, putting $a=4q+1$, $b=4r$ (see Article 14),
\[
(00)=qq-q+rr-1
\]
Since $qq-q$ is evidently always even, $(00)$ will be odd or even according as $r$ is even or odd. Thus $2$ will belong to $A$ or to $C$ according as $b$ has the form $8m$ or the form $8m+4$. This is precisely the theorem found by inductive investigation in Article 14.

II. We can treat the other case, where $p$ has the form $8n+5$, just as completely. Here the number $2$ belongs either to $B$ or to $D$. We readily see that in the former case $\frac12(p-1)$ belongs to $B$ and $\frac12(p+1)$ to $D$, whereas in the latter $\frac12(p-1)$ belongs to $D$ and $\frac12(p+1)$ to $B$. Now consider that, if $\beta$ is a number from $B$ followed by a number from $D$, then $p-\beta-1$ will also be a number from $B$ and $p-\beta$ from $D$; that is, numbers with this property always occur in pairs. Their number, $(13)$, will therefore be even, except when one of them is paired with itself, that is, when $\frac12(p-1)$ belongs to $B$ and $\frac12(p+1)$ to $D$; in this case $(13)$ will be odd. From this we conclude that $(13)$ is even whenever $2$ belongs to $D$, and odd whenever $2$ belongs to $B$. But we have
\[
16(13)=aa+bb+2a+4b+1
\]

\SourceLeaf{99}%
or, putting $a=4q+1$, $b=4r+2$,
\[
(13)=qq+q+rr+2r+1
\]
Thus $(13)$ will be odd whenever $r$ is even; conversely, $(13)$ will be even whenever $r$ is odd. From this we conclude that $2$ belongs to $B$ whenever $b$ has the form $8m+2$, and to $D$ whenever $b$ has the form $8m+6$.

The result of these investigations can be stated as follows:

\emph{The number $2$ belongs to the collection $A,B,C$ or $D$ according as the number $\frac12b$ has the form $4m$, $4m+1$, $4m+2$ or $4m+3$.}

\Article{22}
In the \emph{Disquisitiones Arithmeticae} we explained the general theory of circle division and of the solution of the equation $x^p-1=0$. Among other things, we showed that, if $\mu$ is a divisor of the number $p-1$, the function $\frac{x^p-1}{x-1}$ can be split into $\mu$ factors of degree $\frac{p-1}{\mu}$ with the help of an auxiliary equation of degree $\mu$. Besides the general theory of this factorization, we separately considered the special cases where $\mu=2$ or $\mu=3$ in Articles 356--358 of that work, and showed how to give the auxiliary equation a priori, that is, without working out the scheme of least residues of the powers of a primitive root for the modulus $p$. Even without our pointing it out, attentive readers will readily see the very close connection between the next case of that theory, namely $\mu=4$, and the investigations developed here in Articles 15--20, with whose help that case too can be treated completely without difficulty. But we reserve this treatment for another occasion, and accordingly have preferred to carry out the investigation in this paper in a purely arithmetical form, without bringing in the theory of the equation $x^p-1=0$ at all. Instead, as a finishing touch, we shall add some further new, purely arithmetical theorems very closely connected with the subject treated so far.
